%============================================================
%  Bd. 12 - Halbordnungen
%============================================================

\documentclass[main.tex]{subfiles}

\ifSubfilesClassLoaded{
  \BandSetupStandalone{B12}
}{
}

\title{Bd. 12 - Halbordnungen}
\author{Martin Kunze}
\date{}
\setcounter{file}{12}

\hypersetup{
  pdftitle={Bd. 12 - Halbordnungen},
  pdfauthor={Martin Kunze}
}

\begin{document}

\title{Bd. 12 - Halbordnungen}
\author{Martin Kunze}
\date{}
\hypersetup{pageanchor=false}
\maketitle
\hypersetup{pageanchor=true}
\setcounter{tocdepth}{3}
\tableofcontents
\setcounter{file}{12}
\renewcommand{\FormulaBandID}{12}
\chapter{Einleitung}

Dieser Band führt Halbordnungen, auch partielle Ordnungen genannt, als
reflexive, transitive und antisymmetrische Relationsoperatoren ein. Er
entwickelt den dualen Ordnungsoperator, minimale und maximale Elemente,
Minimum und Maximum, Restriktionen sowie Hauptfilter. Schranken und
Supremumsbegriffe folgen in Band~13; Paarinfima und Paarsuprema in Band~14,
totale Ordnungen in Band~15 und Wohlordnungen in Band~16.

\chapter{Halbordnungen}

\section{Begriff der Halbordnung}
\subsection{Axiome einer partiellen Ordnung}

\FormulaAxiomDeltaK[Antisymmetrie]{%
  x\leq y \dsep y\leq x \vdash x=y%
}{OrderAntisymmetry}{
  \DeltaRow{Mengen}{A \dsep x \dsep y}
  \DeltaRow{Relationsoperatoren}{\leq}
}

\subsection{Definition der partiellen Ordnung}

\FormulaDefDeltaK[Begriff der partiellen Ordnung]{\PartOrd{A,\leq}}{Partielle Ordnung}{
  \DeltaRow{Mengen}{A \dsep x \dsep y \dsep z}
  \DeltaRow{Relationsoperatoren}{\leq}
  \DeltaRow{\textbf{Axiome}}{}
  \DeltaRow{Binärer Relationsoperator}
           {\leq\ \text{auf }A}
  \DeltaRow{Reflexivität}
           {x\in A \vdash x\leq x}
           [\FormulaRefAuto{EqRelReflexivity}[ax]]
  \DeltaRow{Transitivität}
           {x,y,z\in A \dsep x\leq y \dsep y\leq z \vdash x\leq z}
           [\FormulaRefAuto{EqRelTransitivity}[ax]]
  \DeltaRow{Antisymmetrie}
           {x\leq y \dsep y\leq x \vdash x=y}
           [\FormulaRefAuto{x\leq y \dsep y\leq x \vdash x=y}[ax]]
  \DeltaRow{\textbf{Neue Symbole}}{}
  \DeltaPrem{Partielle Ordnung}{\PartOrd{A,\leq}}
}

Die Ordnung wird somit stets zusammen mit ihrem Träger angegeben. Ihre
Auswertung erfolgt in Infixnotation: \(x\leq y\), nicht durch die
Zugehörigkeit eines geordneten Paares zu einer Relationsmenge.

\subsection{Dualer Relationsoperator einer partiellen Ordnung}

\FormulaDefDeltaK[Dualer Relationsoperator auf einem Träger]{%
  x,y\in A\vdash
  \bigl(x\geq y\leftrightarrow y\leq x\bigr)
}{OrderDualRelation}{%
  \DeltaRow{Mengen}{A\dsep x\dsep y}
  \DeltaRow{Relationsoperatoren}{\leq\dsep\geq}
  \DeltaRow{Auswertung}{x,y\in A\vdash(x\geq y\leftrightarrow y\leq x)}
  \DeltaRow{\textbf{Neues Symbol}}{}
  \DeltaPrem{Dualer Relationsoperator}{\geq}
}

\FormulaThmDeltaK[Der duale Relationsoperator ist eine partielle Ordnung]{%
  \PartOrd{A,\leq}\vdash\PartOrd{A,\geq}
}{OrderDualPartialOrder}{%
  \DeltaRow{Mengen}{A\dsep x\dsep y\dsep z}
  \DeltaRow{Relationsoperatoren}{\leq\dsep\geq}
}
\begin{tabproofsplitwide}
  \proofpartwideindR[Reflexivität]{%
    \PartOrd{A,\leq}\dsep x\in A\vdash x\geq x
  }{%
    \PartOrd{A,\leq}\dsep x\in A\vdash x\geq x
  }[OrderDualReflexive]
    \proofstepwidestar[1]{\PartOrd{A,\leq}}{\rA}
    \proofstepwidestar[2]{x\in A}{\rA}
    \proofstepwidestar[1,2]{x\leq x}{%
      \FormulaRefAuto{EqRelReflexivity}[ax]{2}}
    \proofstepwidestar[2]{x\geq x\leftrightarrow x\leq x}{%
      \FormulaRefAuto{OrderDualRelation}[def]{2,2}}
    \proofstepwidestar[1,2]{x\geq x}{%
      \FormulaRefAuto{P\leftrightarrow Q, Q\vdash P}{4,3}}
  \closeproofpartwideindR

  \proofpartwideindR[Transitivität]{%
    \begin{aligned}[t]
      &\PartOrd{A,\leq}\dsep x,y,z\in A
        \dsep x\geq y\dsep y\geq z\vdash{}\\
      &x\geq z
    \end{aligned}
  }{%
    \begin{aligned}[t]
      &\PartOrd{A,\leq}\dsep x,y,z\in A
        \dsep x\geq y\dsep y\geq z\vdash{}\\
      &x\geq z
    \end{aligned}
  }[OrderDualTransitive]
    \proofstepwidestar[1]{\PartOrd{A,\leq}}{\rA}
    \proofstepwidestar[2]{x\in A}{\rA}
    \proofstepwidestar[3]{y\in A}{\rA}
    \proofstepwidestar[4]{z\in A}{\rA}
    \proofstepwidestar[5]{x\geq y}{\rA}
    \proofstepwidestar[6]{y\geq z}{\rA}
    \proofstepwidestar[2,3,5]{y\leq x}{%
      \FormulaRefAuto{P\leftrightarrow Q, P\vdash Q}{%
        \FormulaRefAuto{OrderDualRelation}[def]{2,3},5}}
    \proofstepwidestar[3,4,6]{z\leq y}{%
      \FormulaRefAuto{P\leftrightarrow Q, P\vdash Q}{%
        \FormulaRefAuto{OrderDualRelation}[def]{3,4},6}}
    \proofstepwidestar[1,2,3,4,5,6]{z\leq x}{%
      \FormulaRefAuto{EqRelTransitivity}[ax]{4,3,2,8,7}}
    \proofstepwidestar[2,4]{x\geq z\leftrightarrow z\leq x}{%
      \FormulaRefAuto{OrderDualRelation}[def]{2,4}}
    \proofstepwidestar[1,2,3,4,5,6]{x\geq z}{%
      \FormulaRefAuto{P\leftrightarrow Q, Q\vdash P}{10,9}}
  \closeproofpartwideindR

  \proofpartwideindR[Antisymmetrie]{%
    \begin{aligned}[t]
      &\PartOrd{A,\leq}\dsep x,y\in A
        \dsep x\geq y\dsep y\geq x\vdash x=y
    \end{aligned}
  }{%
    \begin{aligned}[t]
      &\PartOrd{A,\leq}\dsep x,y\in A
        \dsep x\geq y\dsep y\geq x\vdash x=y
    \end{aligned}
  }[OrderDualAntisymmetric]
    \proofstepwidestar[1]{\PartOrd{A,\leq}}{\rA}
    \proofstepwidestar[2]{x\in A}{\rA}
    \proofstepwidestar[3]{y\in A}{\rA}
    \proofstepwidestar[4]{x\geq y}{\rA}
    \proofstepwidestar[5]{y\geq x}{\rA}
    \proofstepwidestar[2,3,4]{y\leq x}{%
      \FormulaRefAuto{P\leftrightarrow Q, P\vdash Q}{%
        \FormulaRefAuto{OrderDualRelation}[def]{2,3},4}}
    \proofstepwidestar[2,3,5]{x\leq y}{%
      \FormulaRefAuto{P\leftrightarrow Q, P\vdash Q}{%
        \FormulaRefAuto{OrderDualRelation}[def]{3,2},5}}
    \proofstepwidestar[1,2,3,4,5]{x=y}{%
      \FormulaRefAuto{x\leq y \dsep y\leq x\vdash x=y}[ax]{7,6}}
  \closeproofpartwideindR

  \proofpartwide{Schluss}
    \proofstepwidestar[1]{\PartOrd{A,\leq}}{\rA}
    \proofstepwidestar[2]{x\in A}{\rA}
    \proofstepwidestar[1,2]{x\geq x}{%
      \FormulaRefAuto{OrderDualReflexive}[thm]{1,2}}
    \proofstepwidestar[4]{y\in A}{\rA}
    \proofstepwidestar[5]{z\in A}{\rA}
    \proofstepwidestar[6]{x\geq y}{\rA}
    \proofstepwidestar[7]{y\geq z}{\rA}
    \proofstepwidestar[1,2,4,5,6,7]{x\geq z}{%
      \FormulaRefAuto{OrderDualTransitive}[thm]{1,2,4,5,6,7}}
    \proofstepwidestar[9]{y\geq x}{\rA}
    \proofstepwidestar[1,2,4,6,9]{x=y}{%
      \FormulaRefAuto{OrderDualAntisymmetric}[thm]{1,2,4,6,9}}
    \proofstepwidestar[1]{\PartOrd{A,\geq}}{%
      \FormulaRefAuto{Partielle Ordnung}[def]{3,8,10}}
  \closeproofpartwide
\end{tabproofsplitwide}


\section{Extremale Elemente}

\subsection{Minimales Element und Minimum}

Ein minimales Element besitzt in der Teilmenge kein echt kleineres Element.
Ein Minimum ist demgegenüber ein kleinstes Element und liegt unter
\emph{jedem} Element der Teilmenge. Die folgenden Sätze halten den genauen
Zusammenhang fest.

\FormulaDefDeltaK[Minimales Element einer Teilmenge]{%
  T\subseteq A\vdash
  \operatorname{MinEl}_{A,\leq}(m,T)
  \coloneqq
  m\in T\land\forall x\in T\,(x\leq m\rightarrow x=m)
}{OrderMinimalElement}{%
  \DeltaRow{Mengen}{A\dsep T\dsep m\dsep x}
  \DeltaRow{Relationsoperatoren}{\leq}
  \DeltaPrem{Partielle Ordnung}{\PartOrd{A,\leq}}
  \DeltaRow{Teilmenge}{T\subseteq A}
  \DeltaRow{\textbf{Neues Symbol}}{}
  \DeltaPrem{Minimales Element}{\operatorname{MinEl}_{A,\leq}(m,T)}
}

% ------------------------------------------------------------
% (1) Eindeutigkeit des Minimum-Kriteriums
% ------------------------------------------------------------
\FormulaThmDeltaK[Eindeutigkeit eines Minimums]{%
  \exists m\in T\,\forall x\in T\,(m\leq x)
  \vdash \exists! m\in T\,\forall x\in T\,(m\leq x)%
}{MinimumUnique}{
  \DeltaRow{Mengen}{A \dsep T \dsep m \dsep n \dsep x}
  \DeltaRow{Relationsoperatoren}{\leq}
  \DeltaPrem{Partielle Ordnung}{\PartOrd{A,\leq}}
}
\begin{tabproof}
  \proofstep{1}{\exists m\in T\,\forall x\in T\,(m\leq x)}{\rA}

  \proofstep{2}{m\in T \land \forall x\in T\,(m\leq x)}{\rA}
  \proofstep{3}{n\in T \land \forall x\in T\,(n\leq x)}{\rA}

  \proofstep{2}{m\in T}{\rAEa{2}}
  \proofstep{2}{\forall x\in T\,(m\leq x)}{\rAEb{2}}
  \proofstep{3}{n\in T}{\rAEa{3}}
  \proofstep{3}{\forall x\in T\,(n\leq x)}{\rAEb{3}}

  \proofstep{2,3}{m\leq n}{\rRE{6,\rUE{5}}}
  \proofstep{2,3}{n\leq m}{\rRE{4,\rUE{7}}}

  \proofstep{2,3}{m=n}{%
    \FormulaRefAuto{x\leq y \dsep y\leq x\vdash x=y}[ax]{8,9}}

  \proofstep{1}{\exists! m\in T\,\forall x\in T\,(m\leq x)}{%
    \UEI{1,2,3,10}}
\end{tabproof}


% ------------------------------------------------------------
% (2) Definition: Minimum als \iota-Term (partielle Definition)
% ------------------------------------------------------------
\FormulaDefDeltaK[Minimum]{%
  \begin{aligned}
    &\exists m\in T\,\forall x\in T\,(m\leq x)\\
    &\vdash \min_{A,\leq}(T)\coloneqq
      \iota m\in T\,\forall x\in T\,(m\leq x)
  \end{aligned}%
}{Minimum}{
  \DeltaRow{Mengen}{A \dsep T \dsep m \dsep x}
  \DeltaRow{Relationsoperatoren}{\leq}
  \DeltaPrem{Partielle Ordnung}{\PartOrd{A,\leq}}
}

\FormulaThmDeltaK[Ein Minimum ist ein minimales Element]{%
  \begin{aligned}[t]
    &T\subseteq A\dsep
      \exists q\in T\,\forall x\in T\,(q\leq x)\vdash{}\\
    &\operatorname{MinEl}_{A,\leq}\bigl(\min_{A,\leq}(T),T\bigr)
  \end{aligned}
}{MinimumIsMinimalElement}{%
  \DeltaRow{Mengen}{A\dsep T\dsep q\dsep x}
  \DeltaRow{Relationsoperatoren}{\leq}
  \DeltaPrem{Partielle Ordnung}{\PartOrd{A,\leq}}
}
\begin{tabproofwide}
  \proofstepwidestar[1]{T\subseteq A}{\rA}
  \proofstepwidestar[2]{%
    \exists q\in T\,\forall x\in T\,(q\leq x)}{\rA}
  \proofstepwidestar[2]{\min_{A,\leq}(T)\in T}{%
    \FormulaRefAuto{Minimum}[def]{2}}
  \proofstepwidestar[4]{x\in T}{\rA}
  \proofstepwidestar[5]{x\leq\min_{A,\leq}(T)}{\rA}
  \proofstepwidestar[2,4]{\min_{A,\leq}(T)\leq x}{%
    \FormulaRefAuto{Minimum}[def]{2,4}}
  \proofstepwidestar[2,4,5]{x=\min_{A,\leq}(T)}{%
    \FormulaRefAuto{x\leq y\dsep y\leq x\vdash x=y}[ax]{5,6}}
  \proofstepwidestar[2,4]{%
    x\leq\min_{A,\leq}(T)\rightarrow x=\min_{A,\leq}(T)}{%
    \rRI{5,7}}
  \proofstepwidestar[2]{%
    \forall x\in T\,
      \bigl(x\leq\min_{A,\leq}(T)\rightarrow x=\min_{A,\leq}(T)\bigr)}{%
    \rUI{\rRI{4,8}}}
  \proofstepwidestar[1,2]{%
    \operatorname{MinEl}_{A,\leq}\bigl(\min_{A,\leq}(T),T\bigr)}{%
    \FormulaRefAuto{OrderMinimalElement}[def]{1,\rAI{3,9}}}
\end{tabproofwide}

\subsection{Maximales Element und Maximum}

\FormulaDefDeltaK[Maximales Element einer Teilmenge]{%
  T\subseteq A\vdash
  \operatorname{MaxEl}_{A,\leq}(m,T)
  \coloneqq
  m\in T\land\forall x\in T\,(m\leq x\rightarrow x=m)
}{OrderMaximalElement}{%
  \DeltaRow{Mengen}{A\dsep T\dsep m\dsep x}
  \DeltaRow{Relationsoperatoren}{\leq}
  \DeltaPrem{Partielle Ordnung}{\PartOrd{A,\leq}}
  \DeltaRow{Teilmenge}{T\subseteq A}
  \DeltaRow{\textbf{Neues Symbol}}{}
  \DeltaPrem{Maximales Element}{\operatorname{MaxEl}_{A,\leq}(m,T)}
}

% ------------------------------------------------------------
% (3) Eindeutigkeit des Maximum-Kriteriums (kurz mit  \exists m\in T ...)
% ------------------------------------------------------------
\FormulaThmDeltaK[Eindeutigkeit eines Maximums]{%
  \exists m\in T\,\forall x\in T\,(x\leq m)
  \vdash \exists! m\in T\,\forall x\in T\,(x\leq m)%
}{MaximumUnique}{
  \DeltaRow{Mengen}{A \dsep T \dsep m \dsep n \dsep x}
  \DeltaRow{Relationsoperatoren}{\leq}
  \DeltaPrem{Partielle Ordnung}{\PartOrd{A,\leq}}
}
\begin{tabproof}
  \proofstep{1}{\exists m\in T\,\forall x\in T\,(x\leq m)}{\rA}

  \proofstep{2}{m\in T \land \forall x\in T\,(x\leq m)}{\rA}
  \proofstep{3}{n\in T \land \forall x\in T\,(x\leq n)}{\rA}

  \proofstep{2}{m\in T}{\rAEa{2}}
  \proofstep{2}{\forall x\in T\,(x\leq m)}{\rAEb{2}}
  \proofstep{3}{n\in T}{\rAEa{3}}
  \proofstep{3}{\forall x\in T\,(x\leq n)}{\rAEb{3}}

  \proofstep{2,3}{m\leq n}{\rRE{4,\rUE{7}}}
  \proofstep{2,3}{n\leq m}{\rRE{6,\rUE{5}}}

  \proofstep{2,3}{m=n}{%
    \FormulaRefAuto{x\leq y \dsep y\leq x\vdash x=y}[ax]{8,9}}

  \proofstep{1}{\exists! m\in T\,\forall x\in T\,(x\leq m)}{%
    \UEI{1,2,3,10}}
\end{tabproof}

% ------------------------------------------------------------
% (4) Definition: Maximum als \iota-Term (partielle Definition)
% ------------------------------------------------------------
\FormulaDefDeltaK[Maximum]{%
  \begin{aligned}
    &\exists m\in T\,\forall x\in T\,(x\leq m)\\
    &\vdash \max_{A,\leq}(T)\coloneqq
      \iota m\in T\,\forall x\in T\,(x\leq m)
  \end{aligned}%
}{Maximum}{
  \DeltaRow{Mengen}{A \dsep T \dsep m \dsep x}
  \DeltaRow{Relationsoperatoren}{\leq}
  \DeltaPrem{Partielle Ordnung}{\PartOrd{A,\leq}}
}

\FormulaThmDeltaK[Ein Maximum ist ein maximales Element]{%
  \begin{aligned}[t]
    &T\subseteq A\dsep
      \exists q\in T\,\forall x\in T\,(x\leq q)\vdash{}\\
    &\operatorname{MaxEl}_{A,\leq}\bigl(\max_{A,\leq}(T),T\bigr)
  \end{aligned}
}{MaximumIsMaximalElement}{%
  \DeltaRow{Mengen}{A\dsep T\dsep q\dsep x}
  \DeltaRow{Relationsoperatoren}{\leq}
  \DeltaPrem{Partielle Ordnung}{\PartOrd{A,\leq}}
}
\begin{tabproofwide}
  \proofstepwidestar[1]{T\subseteq A}{\rA}
  \proofstepwidestar[2]{%
    \exists q\in T\,\forall x\in T\,(x\leq q)}{\rA}
  \proofstepwidestar[2]{\max_{A,\leq}(T)\in T}{%
    \FormulaRefAuto{Maximum}[def]{2}}
  \proofstepwidestar[4]{x\in T}{\rA}
  \proofstepwidestar[5]{\max_{A,\leq}(T)\leq x}{\rA}
  \proofstepwidestar[2,4]{x\leq\max_{A,\leq}(T)}{%
    \FormulaRefAuto{Maximum}[def]{2,4}}
  \proofstepwidestar[2,4,5]{x=\max_{A,\leq}(T)}{%
    \FormulaRefAuto{x\leq y\dsep y\leq x\vdash x=y}[ax]{6,5}}
  \proofstepwidestar[2,4]{%
    \max_{A,\leq}(T)\leq x\rightarrow x=\max_{A,\leq}(T)}{%
    \rRI{5,7}}
  \proofstepwidestar[2]{%
    \forall x\in T\,
      \bigl(\max_{A,\leq}(T)\leq x\rightarrow x=\max_{A,\leq}(T)\bigr)}{%
    \rUI{\rRI{4,8}}}
  \proofstepwidestar[1,2]{%
    \operatorname{MaxEl}_{A,\leq}\bigl(\max_{A,\leq}(T),T\bigr)}{%
    \FormulaRefAuto{OrderMaximalElement}[def]{1,\rAI{3,9}}}
\end{tabproofwide}


\section{Restriktion einer partiellen Ordnung}

\FormulaDefDeltaK[Restriktion eines Ordnungsoperators]{%
  x,y\in T\vdash
  \bigl(x\leq_T y\leftrightarrow x\leq y\bigr)
}{OrderRestrictionOperator}{%
  \DeltaRow{Mengen}{A\dsep T\dsep x\dsep y}
  \DeltaRow{Relationsoperatoren}{\leq\dsep\leq_T}
  \DeltaPrem{Teilmenge}{T\subseteq A}
  \DeltaRow{\textbf{Neues Symbol}}{}
  \DeltaPrem{Restriktionsoperator}{\leq_T}
}

Der Index \(T\) macht den neuen Träger sichtbar. Die Restriktion ändert
also nicht die Vergleichsaussage, sondern nur den Bereich, auf dem der
Relationsoperator ausgewertet wird.

\FormulaThmDeltaK[Reflexivität der Restriktion]{%
  T\subseteq A\dsep x\in T\vdash x\leq_T x
}{OrderRestrictionReflexive}{%
  \DeltaRow{Mengen}{A\dsep T\dsep x}
  \DeltaRow{Relationsoperatoren}{\leq\dsep\leq_T}
  \DeltaPrem{Partielle Ordnung}{\PartOrd{A,\leq}}
}
\begin{tabproof}
  \proofstep{1}{T\subseteq A}{\rA}
  \proofstep{2}{x\in T}{\rA}
  \proofstep{1,2}{x\in A}{%
    \FormulaRefAuto{A\subseteq B, x\in A \vdash x\in B}{1,2}}
  \proofstep{1,2}{x\leq x}{%
    \FormulaRefAuto{EqRelReflexivity}[ax]{3}}
  \proofstep{2}{x\leq_T x\leftrightarrow x\leq x}{%
    \FormulaRefAuto{OrderRestrictionOperator}[def]{2,2}}
  \proofstep{1,2}{x\leq_T x}{%
    \FormulaRefAuto{P\leftrightarrow Q, Q\vdash P}{5,4}}
\end{tabproof}

\FormulaThmDeltaK[Transitivität der Restriktion]{%
  x,y,z\in T\dsep x\leq_T y\dsep y\leq_T z\vdash x\leq_T z
}{OrderRestrictionTransitive}{%
  \DeltaRow{Mengen}{A\dsep T\dsep x\dsep y\dsep z}
  \DeltaRow{Relationsoperatoren}{\leq\dsep\leq_T}
  \DeltaPrem{Teilmenge}{T\subseteq A}
  \DeltaPrem{Partielle Ordnung}{\PartOrd{A,\leq}}
}
\begin{tabproofwide}
  \proofstepwidestar[1]{x\in T}{\rA}
  \proofstepwidestar[2]{y\in T}{\rA}
  \proofstepwidestar[3]{z\in T}{\rA}
  \proofstepwidestar[4]{x\leq_T y}{\rA}
  \proofstepwidestar[5]{y\leq_T z}{\rA}
  \proofstepwidestar[1,2,3]{x,y,z\in A}{%
    \rAI{%
      \FormulaRefAuto{A\subseteq B, x\in A \vdash x\in B}{T\subseteq A,1},%
      \rAI{%
        \FormulaRefAuto{A\subseteq B, x\in A \vdash x\in B}{T\subseteq A,2},%
        \FormulaRefAuto{A\subseteq B, x\in A \vdash x\in B}{T\subseteq A,3}}}}
  \proofstepwidestar[1,2,4]{x\leq y}{%
    \FormulaRefAuto{P\leftrightarrow Q, P\vdash Q}{%
      \FormulaRefAuto{OrderRestrictionOperator}[def]{1,2},4}}
  \proofstepwidestar[2,3,5]{y\leq z}{%
    \FormulaRefAuto{P\leftrightarrow Q, P\vdash Q}{%
      \FormulaRefAuto{OrderRestrictionOperator}[def]{2,3},5}}
  \proofstepwidestar[1,2,3,4,5]{x\leq z}{%
    \FormulaRefAuto{EqRelTransitivity}[ax]{6,7,8}}
  \proofstepwidestar[1,3]{x\leq_T z\leftrightarrow x\leq z}{%
    \FormulaRefAuto{OrderRestrictionOperator}[def]{1,3}}
  \proofstepwidestar[1,2,3,4,5]{x\leq_T z}{%
    \FormulaRefAuto{P\leftrightarrow Q, Q\vdash P}{10,9}}
\end{tabproofwide}


\FormulaThmDeltaK[Antisymmetrie der Restriktion]{%
  x,y\in T\dsep x\leq_T y\dsep y\leq_T x\vdash x=y
}{OrderRestrictionAntisymmetric}{%
  \DeltaRow{Mengen}{A\dsep T\dsep x\dsep y}
  \DeltaRow{Relationsoperatoren}{\leq\dsep\leq_T}
  \DeltaPrem{Teilmenge}{T\subseteq A}
  \DeltaPrem{Partielle Ordnung}{\PartOrd{A,\leq}}
}
\begin{tabproofwide}
  \proofstepwidestar[1]{x\in T}{\rA}
  \proofstepwidestar[2]{y\in T}{\rA}
  \proofstepwidestar[3]{x\leq_T y}{\rA}
  \proofstepwidestar[4]{y\leq_T x}{\rA}
  \proofstepwidestar[1,2,3]{x\leq y}{%
    \FormulaRefAuto{P\leftrightarrow Q, P\vdash Q}{%
      \FormulaRefAuto{OrderRestrictionOperator}[def]{1,2},3}}
  \proofstepwidestar[1,2,4]{y\leq x}{%
    \FormulaRefAuto{P\leftrightarrow Q, P\vdash Q}{%
      \FormulaRefAuto{OrderRestrictionOperator}[def]{2,1},4}}
  \proofstepwidestar[1,2,3,4]{x=y}{%
    \FormulaRefAuto{x\leq y \dsep y\leq x \vdash x=y}[ax]{5,6}}
\end{tabproofwide}


\FormulaThmDeltaK[Restriktion einer partiellen Ordnung]{%
  T\subseteq A\dsep\PartOrd{A,\leq}\vdash\PartOrd{T,\leq_T}
}{OrderRestrictionPartialOrder}{%
  \DeltaRow{Mengen}{A\dsep T\dsep x\dsep y\dsep z}
  \DeltaRow{Relationsoperatoren}{\leq\dsep\leq_T}
  \DeltaPrem{Teilmenge}{T\subseteq A}
  \DeltaPrem{Partielle Ordnung}{\PartOrd{A,\leq}}
}
\begin{tabproofwide}
  \proofstepwidestar[1]{T\subseteq A}{\rA}
  \proofstepwidestar[2]{\PartOrd{A,\leq}}{\rA}
  \proofstepwidestar[3]{x\in T}{\rA}
  \proofstepwidestar[1,2,3]{x\leq_T x}{%
    \FormulaRefAuto{OrderRestrictionReflexive}[thm]{1,3}}
  \proofstepwidestar[5]{y\in T}{\rA}
  \proofstepwidestar[6]{z\in T}{\rA}
  \proofstepwidestar[7]{x\leq_T y}{\rA}
  \proofstepwidestar[8]{y\leq_T z}{\rA}
  \proofstepwidestar[1,2,3,5,6,7,8]{x\leq_T z}{%
    \FormulaRefAuto{OrderRestrictionTransitive}[thm]{3,5,6,7,8}}
  \proofstepwidestar[10]{y\leq_T x}{\rA}
  \proofstepwidestar[1,2,3,5,7,10]{x=y}{%
    \FormulaRefAuto{OrderRestrictionAntisymmetric}[thm]{3,5,7,10}}
  \proofstepwidestar[1,2]{\PartOrd{T,\leq_T}}{%
    \FormulaRefAuto{Partielle Ordnung}[def]{4,9,11}}
\end{tabproofwide}

\chapter{Ordnungsbezogene Teilmengen}

\section{Hauptfilter}

\FormulaDefDeltaK[Hauptfilter eines Elements]{%
  \PrincipalFilter{A,\leq}{j}
  \coloneqq \{x\in A\mid j\leq x\}%
}{PrincipalFilter}{%
  \DeltaRow{Mengen}{A\dsep j\dsep x}
  \DeltaRow{Relationsoperatoren}{\leq}
  \DeltaRow{\textbf{Neues Symbol}}{}
  \DeltaPrem{Hauptfilter}{\PrincipalFilter{A,\leq}{j}}
}
\leavevmode\par\medskip

\FormulaThmDeltaK[Hauptfilter liegen im Ordnungsträger]{%
  \PrincipalFilter{A,\leq}{j}\subseteq A
}{PrincipalFilterSubsetCarrier}{%
  \DeltaRow{Mengen}{A\dsep j\dsep x}
  \DeltaRow{Relationsoperatoren}{\leq}
}
\begin{tabproof}
  \proofstep{1}{x\in\PrincipalFilter{A,\leq}{j}}{\rA}
  \proofstep{1}{x\in A\land j\leq x}{%
    \FormulaRefAuto{PrincipalFilter}[def]{1}}
  \proofstep{1}{x\in A}{\rAEa{2}}
  \proofstep{}{\PrincipalFilter{A,\leq}{j}\subseteq A}{%
    \FormulaRefAuto{A \subseteq B \coloneq
      \forall x\,(x\in A\rightarrow x\in B)}{\rUI{\rRI{1,3}}}}
\end{tabproof}

\end{document}
